
\documentclass[oneside,a4paper]{book}
%\pagestyle{headings}


%=============================================================================

\usepackage{amsthm}
\usepackage{xspace}
\usepackage{float}
\usepackage{ifthen}
\usepackage{amsbsy}
\usepackage{amssymb}
\usepackage{balance}
\usepackage{booktabs}
\usepackage{graphicx}
\usepackage{rotating}
\usepackage{multirow}
\usepackage{needspace}
\usepackage{microtype}
\usepackage{bold-extra}
\usepackage{geometry}
\usepackage{varioref}
\usepackage{xcolor}
\usepackage{textcomp}
\usepackage{listings}
\usepackage[normalem]{ulem} %emphasize still italic
\usepackage{ucs}

% \usepackage[utf8]{inputenc}
% \usepackage[htt]{hyphenat}
\usepackage{times}
\usepackage{url}
\usepackage{alltt}
\usepackage{amsmath}
\usepackage{xfrac}
\usepackage{subfigure}
\usepackage{appendix}
\usepackage{stmaryrd}   % for the \shortuparrow
\usepackage[utopia]{quotchap}

\usepackage{setspace}
\usepackage[numbers, sort&compress]{natbib}
\usepackage{mdwlist}        % support for better spaced lists
% allows for temporary adjustment of side margins
\usepackage{chngpage}
\usepackage[normalem]{ulem} 

% constants

\newcounter{qcounter}

% commands
\newcommand{\n}{$\cdot$}
\newcommand{\y}{\checkmark}
\newcommand{\subscript}[1]{$_{\textrm{\footnotesize{#1}}}$}
\newcommand{\superscript}[1]{$^{\textrm{\footnotesize{#1}}}$}
\newcommand{\vertical}[1]{\raisebox{-4em}{\begin{sideways}{#1}\end{sideways}}}

\newboolean{showedits}
\setboolean{showedits}{true} % toggle to show or hide edits
\ifthenelse{\boolean{showedits}}
{
       \newcommand{\ugh}[1]{\textcolor{red}{\uwave{#1}}} % please rephrase
       \newcommand{\ins}[1]{\textcolor{blue}{\uline{#1}}} % please insert
       \newcommand{\del}[1]{\textcolor{red}{\sout{#1}}} % please delete
       \newcommand{\chg}[2]{\textcolor{red}{\sout{#1}}{\ra}\textcolor{blue}{\uline{#2}}} % please change
}{
       \newcommand{\ugh}[1]{#1} % please rephrase
       \newcommand{\ins}[1]{#1} % please insert
       \newcommand{\del}[1]{} % please delete
       \newcommand{\chg}[2]{#2}
}


% ============================================================================
% Put edit comments in a really ugly standout display

\usepackage{xcolor}
\usepackage[normalem]{ulem}
\newcommand{\ra}{$\rightarrow$}


% comments \nb{label}{color}{text}
\newboolean{showcomments}
\setboolean{showcomments}{true}
\ifthenelse{\boolean{showcomments}}
    {\newcommand{\nb}[3]{
        {\colorbox{#2}{\bfseries\sffamily\scriptsize\textcolor{white}{#1}}}
        {\textcolor{#2}{\sf\small$\blacktriangleright$\textit{#3}$\blacktriangleleft$}}}
     \newcommand{\version}{\emph{\scriptsize$-$Id$-$}}
%	 \newcommand{\ugh}[1]{\textcolor{red}{\uwave{#1}}} % please rephrase
%	 \newcommand{\ins}[1]{\textcolor{blue}{\uline{#1}}} % please insert
%	 \newcommand{\del}[1]{\textcolor{red}{\sout{#1}}} % please delete
%	 \newcommand{\chg}[2]{\textcolor{red}{\sout{#1}}{\ra}\textcolor{blue}{\uline{#2}}} % please change
	 \newcommand{\chk}[1]{\textcolor{ForestGreen}{#1}} % changed, please check
	}
    {\newcommand{\nb}[3]{}
     \newcommand{\version}{}
	\newcommand{\chk}[1]{} % changed, please check
	}

% ============================================================================
% Make quotes be italic
\renewenvironment{quote}
    {\list{}{\rightmargin\leftmargin}%
     \item\relax\begin{it}}
    {\end{it}\endlist}

\newcommand{\ttimes}{\ensuremath{\times}}

%=============================================================================

\newcommand{\needlines}[1]{\Needspace{#1\baselineskip}}

% source code
\usepackage{xcolor}
\usepackage{textcomp}
\usepackage{listings}
\definecolor{source}{gray}{0.9}
\lstset{
	language={},
	% characters
	tabsize=3,
	upquote=true,
	escapechar={!},
	keepspaces=true,
	breaklines=false,
	alsoletter={:},
	breakautoindent=true,
	columns=fullflexible,
	showstringspaces=false,
	basicstyle=\footnotesize\ttfamily,
	% background
	frame=single,
    framerule=0pt,
	backgroundcolor=\color{source},
	% numbering
	numbersep=5pt,
	numberstyle=\tiny,
	numberfirstline=true,
	% captioning
	captionpos=b,
	numberbychapter=false,
	% formatting (html)
	moredelim=[is][\textbf]{<b>}{</b>},
	moredelim=[is][\textit]{<i>}{</i>},
	moredelim=[is][\uline]{<u>}{</u>}}
\newcommand{\ct}{\lstinline[backgroundcolor=\color{white},basicstyle=\footnotesize\ttfamily]}
\newcommand{\lct}[1]{{\small\tt #1}}


%----------------------------------------------------------------------------
% references
\newcommand{\tabref}[1]{\hyperref[{tab:#1}]{Table~\ref*{tab:#1}}}
\newcommand{\figref}[1]{\hyperref[{fig:#1}]{Figure~\ref*{fig:#1}}}
\newcommand{\secref}[1]{\hyperref[{sec:#1}]{Section~\ref*{sec:#1}}}
\newcommand{\lstref}[1]{\hyperref[{lst:#1}]{Listing~\ref*{lst:#1}}}
\newcommand{\charef}[1]{\hyperref[{cha:#1}]{Chapter~\ref*{cha:#1}}}
%----------------------------------------------------------------------------

% abbreviations
\tracingcolors 4
\setcounter{tocdepth}{3}
\setcounter{secnumdepth}{3}
\newcommand{\ie}{\emph{i.e.,}\xspace}
\newcommand{\eg}{\emph{e.g.,}\xspace}
\newcommand{\etc}{\emph{etc.}\xspace}
\newcommand{\etal}{\emph{et al.}\xspace}


\newcommand{\newevenside}{
	\ifthenelse{\isodd{\thepage}}{\newpage}{
	\newpage
        \phantom{placeholder} % doesn't appear on page
	\thispagestyle{empty} % if want no header/footer
	\newpage
	}
}

\def\stretchfactor{1}
\newcommand{\mychapter}[1]{\setstretch{1}
    \chapter{#1}\setstretch{\stretchfactor}}

%----------------------------------------------------------------------------
\newcommand{\lessSpace}{\vspace{-1em}}
\DeclareGraphicsExtensions{.pdf,.png}
\graphicspath{{images/}}
\newcommand{\fig}[4]{
	\begin{figure}[#1]
		\centering
		\includegraphics[width=#2\textwidth]{#3}
		\lessSpace
		\caption{\label{fig:#3}#4}
	\end{figure}}

% ===========================================================================


\newcommand{\thesistitle}{Title of the Thesis}
\newcommand{\thesisauthor}{Given Name \& Surname}
\newcommand{\thesisleiter}{Name of the Supervisor}
\newcommand{\thesisasst}{Name of the Assistant}
\newcommand{\thesissubtitle}{Thesis subtitle}
\newcommand{\thesisdate}{Month and Year}



% ===========================================================================

\usepackage[ colorlinks=true, urlcolor=black, linkcolor=black,
			citecolor=black, bookmarksnumbered=true, bookmarks=true,
			plainpages=false,
			pdftitle={\thesistitle}, pdfauthor={\thesisauthor},
			pdfsubject={\thesissubtitle}, pdfpagelabels]{hyperref}

\newcommand{\hrref}[2]{\hyperref}
% ===========================================================================
% ===========================================================================


% D O C U M E N T
% % % % % % % % % % % % % % % % % % % % % % % % % % % % % % % % % %
\begin{document}

% T I T L E
% % % % % % % % % % % % % % % % % % % % % % % % % % % % % % % % % %
\begin{titlepage}  
  \begin{center}  
  
  \begin{figure}[t]  
  \vspace*{-2cm}        % to move header logo at the top 
  \center{\includegraphics[scale=0.2]{logos/MSc_quer.png}}
  \vspace{0.4in}     
  \end{figure}

    \thispagestyle{empty}
    
    {\bfseries\Huge \thesistitle \par
    \Large \vspace{0.1in} \thesissubtitle \par}

    \vspace{0.3in} 
    \LARGE{\textbf{Master Thesis} \\}
    \vspace{0.4in}

    {\Large \thesisauthor}
    
    \vspace{0.3in}
    {\Large Home University \par}
%    {\Large Philosophisch-naturwissenschaftlichen Fakult\"{a}t \\
%            der Universit\"{a}t Bern \par}
    \vfill
    {\Large \thesisdate \par}
  

  \vspace{0.9in}
 
  % === Logos ==============================================     
  \begin{figure}[htp]
    \centering
    \includegraphics[scale=0.30]{logos/UNI_Bern.png}\hfill
    \includegraphics[scale=0.30]{logos/UNI_Neuenburg.png}\hfill
    \includegraphics[scale=0.80]{logos/UNI_Fribourg.png}
  \end{figure}
  % === // Logos ===========================================    


  \end{center}

\end{titlepage}


% A B S T R A C T
% % % % % % % % % % % % % % % % % % % % % % % % % % % % % % % % % %
\chapter*{\centering Abstract}
\begin{quotation}
\noindent 
Abstract (max. 1 page)

Name of the Supervisor, Group, Institute, University, Supervisor

Name of the Assistant, Group, Institute, University, Assistant
\end{quotation}
\clearpage


% C O N T E N T S 
% % % % % % % % % % % % % % % % % % % % % % % % % % % % % % % % % % % % % % % %
\tableofcontents


\section{Abstract}

In this thesis we apply different techniques of machine learning to a real life dataset, gathered anonymously from users of an Android app that collects data about the stations the user has passed. Our goal is to best predict the next station that the user is going to, based on information that would be available to us in the moment. Such information could include time, day, current location, past locations.

We show ways to analyze and optimize the raw data, prepare and filter it in a way that yields the best results for further analysis. We compare different machine learning algorithms on our dataset and discuss the advantages and disadvantages that we encountered with each algorithm. We then present the experiments that we did on our dataset and evaluate the results and predictions from the algorithms to draw conclusions to the reliability and the meaningfulness of our approaches. We further discuss what future work should or could be done to take advantage of the techniques that we explored in our thesis.

\newpage
\section{Introduction}

Machine Learning has seen an incredible boost of interest in recent years in the research community as well as in industrial applications. Due to the amount of data being generated and gathered across various disciplines new approaches to data analysis have to be discovered. It is no longer feasible and sufficient to manually skim through the data and draw conclusions from these analysis. The process would usually be too slow and the amount of data too overwhelming to have a reasonable process of analysis. This issue has led to the rising of new machine learning techniques and related fields which fall under the broader term of artificial intelligence (AI). Many different algorithms and techniques have been developed to account for different problems, data sets and procedures.

The main goal of our thesis was to evaluate and compare different approaches and machine learning algorithms on an existing data sets. We also evaluate different types of use cases and combinations of features that are available in our data sets. We compare the results of these evaluations and try to find a combination of data preparations, algorithms and evaluations that yield results successful enough to be able to be used in an actual application. 

We have gathered geographical and temporal data of users of an Android application that is used for looking up public transportation information. With the data set we should be able to create anonymized user profiles, the precision of which always reflecting the amount of data we have on each user. Based on the historical information we have about each user we try to predict the future behavior of the user, based on features such as current station, day of the week, time, past station(s). Our goal is to accurately predict the next station the user will go to and therefore be able to support the user in ordinary day to day tasks. 

An example use case would be to create a personal assistant that "knows" a users daily routine from historical data and without the user having to explicitly state it. Combined with other data sources such as geographical locations (home, work, gym, train stations, shopping centers, etc.), calendar entries or current public transport timetables the assistant could provide information about when to leave the current place in order to get to the next place that is to be targeted in the users routine. 

The part of this thesis in such a use case is to provide the underlying machine learning framework that can be included in such an application. For that we take the available data we have and analyze the necessary steps to get predictions as accurate as possible. The conclusion drawn from this thesis can be integrated in a more comprehensive application that covers different use cases. We will cover this in more detail in the chapter about future work (section~\ref{sec:future_work}).

\section{Related Work}

Here goes the related work.
\section{Sequence Prediction}

Here goes the Sequence description.

\subsection{Task Description}

(we model our problem as a sequence prediction problem)


\subsection{Proposed Solutions}
(Detailed Description of Different Classifiers used, and how they are going to be used)
\section{Experiments and Evaluation}
\label{sec:experiments}

Here goes the experiments.


\subsection{Introduction / Procedure}

\subsection{Data Set / Data Analysis}
\label{subsec:data_set}

\subsection{Naive Approach}


\subsection{Machine Learning Results}

\subsubsection{Execution Plan}
\subsubsection{Data Preparation / Statistical Analysis}
\subsubsection{Decision Trees (Random Forest is included)}
\subsubsection{Naive Bayes}
\subsubsection{HMM / Neural Networks}

\subsection{Comparison of Results}
\newpage
\section{Conclusion}

Here goes the conclusion.
\newpage
\chapter{Future Work}
\label{cha:future_work}

Even though we have successfully completed our experiments and have, under certain restrictions to our data, gathered a respectable precision with our predictions, there are still things left to be improved. In this final section we list the tasks that could and / or should be taken up following this thesis.

As a logical next step the approaches that we have implemented can be included in the actual Android application where the data is gathered. If carefully and usefully integrated into the application that the user actually benefits but doesn`t consider it as too intruding or too imprecise in the predictions this could lead to a virtuous cycle. The more use the user has the more he uses the feature, in turn providing our system with more data which generally leads to a better prediction for the user. The introduction of the feature would need to carefully monitored, so that it is only active for user with enough gathered data that the system is confident enough of providing a good estimate. Otherwise the prediction might be too shaky, leaving the user with little benefit and generally loosing prospective users in the process. It might be necessary to do some A/B tests to figure out what the target precision for a prediction needs to be in order to display it to the end user. Depending on how much or how often the user likes to see recommendations the precision could be more or less important and the behavior of the application can be adapted.

In combination with creating an actual integration a more flexible approach to the statistical analysis, the data preparation as well as the options of the machine learning algorithms might be taken. After our analysis phase and some trials on the data we have imposed fixed limits on what conditions our data needs to have in order to be relevant for our analysis. Depending on how the user actually behaves it might make sense to either lift or  further limit certain restrictions, based on the complexity that the users data profile inhibits. The overall precision of the system might not be improved too much, however for certain edge cases this could lead to significant gains or could successfully delay the integration of the system until enough data and information is gathered to be precise enough. 

Another type of prediction that we could try to establish is the current position of the user. We have in our whole thesis only ever predicted where the user might go next, however it could also be interesting to get a best guess for the current location of the user, based on time, day, previous usages of the app, and so on. This might help in getting a faster coarse location, faster than GPS and cell signal triangulation. If the predictions would be precise it could easily lead to a great contribution for the Android application.

On a technical level there are multiple options to go forward. One would be to include a form of global state for the prediction. We have only ever looked at a users data profile in isolation. It might be interesting to evaluate certain patterns that are valid for all users and take these into account when creating a prediction for a single user. This might especially be helpful if the system encounters a case where previous data from the user doesn`t yield a probable prediction. Since humans tend to be herd animals and follow similar patterns this could be a helpful contribution.

Another technical option might be to switch the machine learning framework for something more recent. A good example to evaluate would be Google`s TensorFlow, an open source framework based on Python that is founded on the years of research that Google has put into machine learning. TensorFlow is based on neural networks and can be scaled from a smartphone to a large-scale datacenter. It might be very interesting to see what advantages can be gained using an absolute state-of-the-art framework that is already being used in applications by billions of users. The exact effect it would have on our problem set is hard to predict, but it could very well lead to an improved predictions while making less restrictions on the data. 
%\chapter {Related Work}
%In which we learn what have other done to address similar problems. For example, the work of Star \cite{Star89}
%
%\chapter{The Problem}
%In which we understand what the problem is in detail.
%
%\chapter {The Solution}
%In which you describe your solution.
%
%\chapter {The Validation}
%In which you show how well the solution works.
%
%\chapter {Conclusion and Future Work}
%In which we step back, have a critical look at the entire work, then conclude, and learn what lies beyond this thesis.



%END Doc
%-------------------------------------------------------

\bibliography{thesis}
\bibliographystyle{plain}

\end{document}
